% Part: computability
% Chapter: computability-theory
% Section: def-functions-self-reference

\documentclass[../../../include/open-logic-section]{subfiles}

\begin{document}

\olfileid{cmp}{thy}{slf}
\olsection{اپنے آپ دے حوالے نال فنکشناں دی تعریف}

عام طور اُتے فنکشناں نوں
اپنے آپ دے حوالے نال
تعریف کر سکنا فائدہ مند
ہوندا اے۔ مثال دے طور
اُتے، جے قابلِ حساب
فنکشن $k$، $l$ تے~$m$
دتے ہون، تاں مقرر نقطے
دا معاون قضیہ دسدا اے
کہ اک جزوی قابلِ حساب
فنکشن~$f$ موجود اے
جیڑا ہر~$y$ لئی ہِٹھلی
مساوات پوری کردا اے:
\[
f(y) \simeq
\begin{cases}
k(y) & \text{جے $l(y) = 0$} \\
f(m(y)) & \text{ہور صورت وچ۔}
\end{cases}
\]
ہور خاص طور اُتے، $f$ لبھن لئی
پہلاں ایہ تعریف کرو:
\[
g(x,y) \simeq
\begin{cases}
k(y) & \text{جے $l(y) = 0$} \\
\cfind{x}(m(y)) & \text{ہور صورت وچ}
\end{cases}
\]
تے فیر مقرر نقطے دا
معاون قضیہ ورت کے اک
اشاریہ~$e$ لبھیا جاندا
اے ایداں کہ
$\cfind{e}(y) = g(e,y)$۔

اک ٹھوس مثال ویکھو۔
«سب توں وڈا مشترک پورا
ونڈن والا عدد» (greatest
common divisor) دا فنکشن
$\fn{gcd}(u,v)$ ایداں
تعریف کیتا جا سکدا اے:
\[
\fn{gcd}(u,v) \simeq
\begin{cases}
v & \text{جے $u = 0$} \\
\fn{gcd}(\fn{mod}(v, u), u) & \text{ہور صورت وچ}
\end{cases}
\]
جتھے $\fn{mod}(v, u)$
توں مراد $v$ نوں~$u$
نال ونڈن مگروں بچیا
باقی اے۔ مقرر نقطے
دے معاون قضیے نوں
لا کے پتا لگدا اے کہ
$\fn{gcd}$ جزوی قابلِ
حساب اے۔ (حقیقت وچ،
$y$ نوں جوڑے
$\tuple{u, v}$ دا کوڈ
من کے ایہنوں اُتلی
صورت وچ لکھ سکدے آں۔)
ایس توں بعد~$u$ اُتے
استقرا نال ثابت ہوندا اے کہ
$\fn{gcd}$ دراصل کُل اے۔

بے شک، اپنے آپ دے
حوالے والیاں تعریفاں
پچھلیاں مثالاں توں
کِتے ودھ پیچیدہ وی
بنا سکدے آں۔ بہتیاں
پروگرامنگ زباناں اک
نہ اک طریقے نال فنکشناں
نوں اپنے آپ دے حوالے
نال تعریف کرن دی اجازت
دیندیاں نیں۔ پر دھیان
دیو: ایہ گل، اوہدے
مقابلے وچ قدرے گھٹ
حیران کن اے کہ اسی
اک فنکشن نوں اوس فنکشن
دا حساب کرن والے الگورتھم
دے \emph{اشاریے} دے
حوالے نال تعریف کر
سکدے آں؛ مقرر نقطے
دا قضیہ اپنے پورے عموم
وچ ایہ دوجی سہولت
دیندا اے۔

\end{document}
